\section{Diskussion}
\label{sec:Diskussion}

In diesem Abschnitt werden die Ergebnisse der Versuche diskutiert und mögliche Fehlerquellen
aufgezeigt.

\subsection{Invertierender Verstärker}
In der \ref{subsec:Auswertung_InvertierenderVerstaerker} wurden Abweichungen von unter $\qty{3.5}{\percent}$ gemessen. 
Außerdem weisen die \autoref{plt:invertierenderVerstaerker1}, \ref{plt:invertierenderVerstaerker2} und \ref{plt:invertierenderVerstaerker3}  einen 
zu erwartenden Verlauf auf und lassen sich gut an eine typische Low Pass Filter Funktion \ref{eq:lowPassFilter} fitten.
Die Messungen bestätigen somit die Theoretischen Erwartungen.

\subsection{Integrator}
\label{sec:Integrator_Diskussion}
Die \autoref{plt:integrator} zeigt, dass die Messdaten nicht exakt einem linearen Verlauf folgen. 
Aufgrund der niedrigen Frequenzen kann dies vermutlich auf Untergrundeffekte und allgemeine Schwierigkeiten bei der Darstellung eines sauberen Signals auf dem Oszilloskop zurückgeführt werden.
Zudem könnten Messunsicherheiten in den Bauteilen zu größeren Abweichungen führen.
Im Bereich von etwa $\qty{20}{\hertz}$ bis $\qty{60}{\hertz}$ zeigt sich jedoch ein Verhalten, das dem theoretischen Verlauf besser entspricht.
Dieser Frequenzbereich ist allerdings zu schmal, um eine aussagekräftige Analyse ausschließlich auf ihn zu stützen.
Die experimentell bestimmte Verstärkung weicht um $\qty{25.41}{\percent}$ vom theoretisch erwarteten Wert ab, was innerhalb der Toleranzen der Bauteileigenschaften und der allgemeinen Messunsicherheiten liegt.

Anhand von \autoref{fig:integrator_rechteck} und 
\autoref{fig:integrator_dreieck} ist qualitativ zu 
erkennen, dass das Ausgangssignal des Integrators, bis auf einen Vorfaktor, dem zeitlichen Integral 
des Eingangssignals entspricht. Aus dem Rechtecksignal ergibt sich ein Dreieckssignal und aus dem
Dreieckssignal ein Parabelsignal. Dies entspricht den Erwartungen gemäß \autoref{eq:integrator}.

\subsection{Differenzierer}
\label{sec:Differenzierer_Diskussion}
Die \autoref{plt:differenzierer} zeigt einen linearen Verlauf für höhere Frequenzen und verliert im unteren Frequenzbereich leicht an Linearität.
Dies lässt sich auf Untergrundeffekte zurückführen, da im niedrigen Frequenzbereich die Verstärkung abnimmt und dadurch Hintergrundstörungen dominanter werden.
Dies führt dazu, dass die Fit-Funktion auf natürliche Weise leicht nach oben verschoben wird.
Zukünftige Analysen könnten daher versuchen, den Untergrund im Vorfeld zu messen und rechnerisch zu subtrahieren.
Eine Abweichung von $\qty{7.51}{\percent}$ vom theoretischen Wert bestätigt insgesamt die theoretische Erwartung.

In \autoref{fig:differenzierer_rechteck} und
\autoref{fig:differenzierer_dreieck} sind die Eingangs-
und Ausgangssignale des Differenzierers dargestellt. Es wird erwartet, dass an den Sprüngen des
Rechtecksignals die Amplitude des Differenzierersignals einen Peak zeigt. Dies ist 
in \autoref{fig:differenzierer_rechteck} zu erkennen.
Es sei angemerkt, dass die Spannungsausschläge gedämpfte Schwingungen im Schaltkreis anregen, sodas
die Peaks nicht ideal sind und zu den beobachteten Schwingungen führen.
Ein ähnliches Verhalten ist auch im Fall des Dreiecksignals zu beobachten. Hierbei springt 
das Ausgangssignal zwischen zwei Werten hin und her, wobei ebenfalls gedämpfte Schwingungen
angeregt werden. Bis auf die gedämpften Schwingungen entspricht das Ausgangssignal des Differenzierers
somit den Erwartungen gemäß \autoref{eq:differenzierer}.

\subsection{Generator}
Die gemessene Frequenz von $f_\text{exp} = \qty{1.67}{\kilo\hertz}$ weicht um $\qty{33.2}{\percent}$ vom theoretischen Wert $f_\text{theo} = \qty{2.5}{\kilo\hertz}$ ab.  
Diese Differenz kann durch Bauteiltoleranzen, insbesondere beim Kondensator, sowie durch nicht-ideale Eigenschaften der Operationsverstärker und Versorgungsspannungen erklärt werden.  
Trotz der Abweichung liegt der Wert in einer realistischen Größenordnung und bestätigt qualitativ das erwartete Oszillationsverhalten.

\subsection{Generator mit variierender Amplitude}
\label{sec:Generator_Diskussion}
Der experimentelle Wert und die theoretische Vorhersage für $\tau$ passen im Rahmen der Unsicherheiten nicht zusammen.
Dies kann jedoch dadurch erklärt werden, dass der Wert von $\abs{\eta}$ durch das Potentiometer eingestellt wird und nicht bekannt 
ist. Der theoretische Wert liegt unterhalb des experimentellen Wertes. Dies deutet darauf hin, 
dass gemäß \autoref{eq:generatorMitVariierenderAmplitude} der Wert von $\abs{\eta}$ kleiner ist als der 
angenommene Wert von $\abs{\eta} = 1$.